\documentclass[12pt,a4paper]{article}

%% ---- Fonts & Math ----
\usepackage[utf8]{inputenc}
\usepackage[T1]{fontenc}
\usepackage{amsmath,amssymb,amsthm}  % Load amsthm BEFORE newtxmath to avoid \openbox clash
\usepackage{mathtools}
\usepackage{newtxtext}          % Times-style text (standard for econ journals)
\usepackage{newtxmath}          % Matching math font
\usepackage{microtype}          % Subtle typographic improvements

%% ---- Layout ----
\usepackage{geometry}
\usepackage{setspace}
\usepackage{enumitem}
\usepackage{booktabs}
\usepackage[dvipsnames]{xcolor}  % Color mixing for hyperref link colors
\usepackage{natbib}

%% ---- Hyperref (load last, before cleveref) ----
\usepackage[
  colorlinks=true,
  linkcolor={blue!70!black},
  citecolor={green!50!black},
  urlcolor={blue!60!black},
  bookmarks=true,
  bookmarksnumbered=true,
  pdfstartview=FitH,
  pdftitle={Covenants Eliminate Hold-Up in Relationship Banking Under Moral Hazard},
  pdfauthor={Without Loss},
]{hyperref}
\usepackage[capitalize,nameinlink]{cleveref}

%% ---- Page geometry ----
\geometry{
  left=1.15in,
  right=1.15in,
  top=1.1in,
  bottom=1.1in,
}
\onehalfspacing

%% ---- Theorem environments ----
%%  Plain style: theorems, propositions, lemmas, corollaries (italic body)
\theoremstyle{plain}
\newtheorem{theorem}{Theorem}[section]
\newtheorem{proposition}[theorem]{Proposition}
\newtheorem{lemma}[theorem]{Lemma}
\newtheorem{corollary}[theorem]{Corollary}

%%  Definition style: definitions, assumptions, examples (upright body)
\theoremstyle{definition}
\newtheorem{definition}[theorem]{Definition}
\newtheorem{assumption}{Assumption}
\newtheorem{example}[theorem]{Example}

%%  Remark style: remarks, conjectures (upright body, lighter heading)
\theoremstyle{remark}
\newtheorem{remark}[theorem]{Remark}
\newtheorem{conjecture}[theorem]{Conjecture}

%% ---- QED symbol ----
\renewcommand{\qedsymbol}{$\blacksquare$}

%% ---- Custom commands ----
\newcommand{\E}{\mathbb{E}}
\newcommand{\R}{\mathbb{R}}
\newcommand{\N}{\mathbb{N}}
\newcommand{\Prob}{\mathbb{P}}
\newcommand{\Var}{\mathrm{Var}}
\newcommand{\Cov}{\mathrm{Cov}}
\DeclareMathOperator*{\argmax}{arg\,max}
\DeclareMathOperator*{\argmin}{arg\,min}
\DeclareMathOperator{\supp}{supp}

%% ---- Title ----
\title{Covenants Eliminate Hold-Up in Relationship Banking Under Moral Hazard}
\author{Without Loss}
\date{\today}

\begin{document}

\maketitle

\begin{abstract}
\noindent Rajan (1992) shows that a bank's informational advantage generates hold-up, distorting borrower effort and making arm's-length finance preferable for high-quality firms. That analysis restricts attention to plain-vanilla debt contracts. We show this restriction is consequential. Within the identical environment---preserving moral hazard, limited liability, non-verifiable states, and the owner's continuation bias---a debt-with-covenant contract that transfers the liquidation decision unconditionally to the bank implements the second-best allocation when monitoring is costless. The contract is renegotiation-proof: the bank's unilateral optimum coincides with the efficient liquidation rule in both states. When monitoring is costly, bank finance induces strictly higher effort than arm's-length finance, and a quality threshold governs source choice. The threshold result survives qualitatively, but the mechanism shifts from renegotiation-induced hold-up to the direct resource cost of monitoring.
\end{abstract}

\medskip
\noindent \textbf{Keywords:} bank finance; hold-up; debt covenants; control rights; moral hazard; renegotiation-proofness

\smallskip
\noindent \textbf{JEL Classification:} D86; G21; G32; D82

\bigskip

\section{Introduction}
\label{sec:intro}

\citet{Rajan1992} identifies a fundamental tension in corporate borrowing. An informed
bank can observe the state of the world at the interim date and enforce efficient
liquidation; but the very informational advantage that enables control also confers
bargaining power, allowing the bank to extract surplus from the firm ex post. This
hold-up distorts the firm's ex ante effort incentives, generating an endogenous cost of
bank finance. The paper then studies how maturity, source mix, and priority can
partially circumscribe the bank's rent-extraction without eliminating its control.

The analysis in \citet{Rajan1992} is conducted entirely within the class of
\emph{plain-vanilla debt contracts}: the firm borrows a fixed amount and promises a
single repayment, with no contractual transfer of control rights. This restriction is
maintained throughout---for both short-term and long-term bank contracts---and is
not derived from primitives of the environment. The paper acknowledges in a footnote
that commitment mechanisms could approach the first best under special circumstances,
but does not pursue this observation within the main analysis.

We show that this restriction is consequential. The environment studied by
\citet{Rajan1992} admits a simple \emph{debt-with-covenant} contract---a standard
debt instrument augmented by a covenant that transfers the liquidation decision to the
bank unconditionally---that implements the constrained optimum (second best) for all
parameter values when monitoring is costless. The contract is renegotiation-proof: the
bank has no credible threat to deviate from the efficient liquidation rule once the state
is reached. When monitoring is costly, the same covenant structure implements the
second best conditional on monitoring, and the firm's choice between bank and
arm's-length finance is governed by a clean comparison of second-best payoffs under
each regime.

Our contribution is threefold.
\begin{enumerate}[nosep]
\item We show that within the full model of \citet{Rajan1992}---preserving moral
  hazard, limited liability, non-verifiability of the state, and the owner's
  continuation bias---a debt-with-covenant contract implements the second best when
  monitoring is costless (\cref{p:covenant_implements}).
\item We prove that the covenant contract is renegotiation-proof: neither party has a
  credible incentive to deviate from the contractually specified liquidation rule at
  the interim date (\cref{l:renegotiation_proof}).
\item We characterise the firm's optimal choice between bank and arm's-length finance
  when monitoring is costly, and show that a quality threshold governs this choice in
  the same qualitative direction as \citet{Rajan1992}, but for a different reason: the
  trade-off is driven by monitoring costs rather than by renegotiation distortions
  (\cref{p:quality_threshold}).
\end{enumerate}

The results do not overturn the central message of \citet{Rajan1992}---bank finance
carries an endogenous cost---but they clarify its source. When the contract space is
expanded to include control-rights covenants, the cost of bank finance is not
renegotiation-induced hold-up but the direct resource cost of monitoring. The
qualitative comparative statics with respect to project quality survive, but the
mechanism is distinct.

\section{Model}
\label{sec:model}

We preserve every feature of the environment in \citet{Rajan1992}. The only
modification is the contract space available to the bank.

\subsection{Primitives}

\paragraph{Dates and agents.}
The economy runs for three dates $t \in \{0,1,2\}$. There is a single
owner-managed firm (henceforth the \emph{owner}), a competitive banking sector, and
a large pool of arm's-length investors.

\paragraph{Project technology.}
\begin{assumption}[Technology]
\label{a:technology}
The project requires a fixed investment $I > 0$ at $t=0$. At $t=1$ a state
$s \in \{G,B\}$ is realised. If the project is continued to $t=2$, the payoff is $X$
with probability one in state $G$, and $X$ with probability $q \in (0,1)$ and zero
with probability $1-q$ in state $B$. If the project is liquidated at $t=1$, it yields
$L$. Parameters satisfy
\begin{equation}
  X > I > L > qX.
  \label{eq:param_order}
\end{equation}
\end{assumption}

The inequality $L > qX$ means liquidation is strictly efficient in state $B$, while
$X > I > L$ means continuation is strictly efficient in state $G$ and the project has
positive net present value when the good state is reached.

\paragraph{Moral hazard.}
\begin{assumption}[Moral hazard]
\label{a:moral_hazard}
The owner chooses effort $e \in [0,\infty)$ at $t=0$, after contracting but before the
state is realised. Effort is privately observed by the owner and is not verifiable by
any court. The cost of effort is $c(e) = e$. The probability of the good state is
$p(e) := \Pr[s = G \mid e]$, where $p$ satisfies:
\begin{enumerate}[nosep,label=(\roman*)]
  \item $p'(e) > 0$ and $p''(e) < 0$ for all $e \geq 0$;
  \item $p'(0) = 1/\varepsilon$ for some small $\varepsilon > 0$, and $p'(\infty) = 0$.
\end{enumerate}
\end{assumption}

\begin{assumption}[Limited liability]
\label{a:limited_liability}
The owner has no wealth at $t=0$. All financing is external. The owner's payoff in
any state is bounded below by zero.
\end{assumption}

\paragraph{Information structure.}
\begin{assumption}[Information]
\label{a:information}
At $t=1$, the state $s$ is observed by the owner and, if the bank has monitored, by
the bank. The state is \emph{not} verifiable by any court. Arm's-length investors
observe no information about the state at $t=1$. Monitoring by the bank costs $c \geq 0$
and is incurred at $t=0$ as part of the lending decision.
\end{assumption}

The non-verifiability of the state is the key friction: contracts cannot be written
directly contingent on $s$. This is identical to the assumption in \citet{Rajan1992}.

\paragraph{Continuation bias.}
\begin{assumption}[Continuation bias]
\label{a:continuation_bias}
Because the owner has limited liability and a residual claim, she strictly prefers
continuation to liquidation in both states whenever she controls the liquidation
decision. Formally, her payoff upon liquidation is zero (the bank recovers $L$), while
her payoff upon continuation is strictly positive in expectation.
\end{assumption}

This assumption, maintained throughout \citet{Rajan1992}, implies that the owner
will never voluntarily liquidate. Any efficient liquidation in state $B$ must be
enforced by the lender.

\subsection{Contract spaces}

\begin{definition}[Plain-vanilla debt contract]
\label{d:plain_vanilla}
A \emph{plain-vanilla debt contract} specifies an amount borrowed $A \geq 0$ at some
date $t_0$ and a required repayment $D \geq A$ at some later date $t_1 > t_0$. The
contract contains no provisions governing the liquidation decision. The owner retains
control over continuation at $t=1$.
\end{definition}

This is the contract space studied in \citet{Rajan1992}. Short-term contracts have
$t_1 - t_0 = 1$; long-term contracts have $t_1 - t_0 = 2$.

\begin{definition}[Debt-with-covenant contract]
\label{d:covenant}
A \emph{debt-with-covenant contract} specifies an amount borrowed $I$ at $t=0$, a
required repayment $D_b \geq I$ at $t=2$, and a \emph{control covenant} that
transfers the liquidation decision unconditionally to the bank at $t=1$. The bank
exercises this right after observing the state (having monitored). No other feature of
the environment is altered.
\end{definition}

The covenant is not state-contingent---it cannot be, since the state is not
verifiable---but it transfers the \emph{decision right} to the informed party. The
bank then exercises this right after privately observing $s$.

\begin{remark}
\label{r:covenant_standard}
Control covenants of this form are standard in commercial lending practice. They
include provisions granting the lender the right to accelerate repayment or to
appoint a receiver upon the occurrence of specified financial conditions. The
theoretical device here is the limiting case in which the bank holds unconditional
control rights over the continuation decision.
\end{remark}

\subsection{Timing}

The timing is identical to \citet{Rajan1992}:
\begin{enumerate}[nosep]
\item $t=0$: The owner approaches lenders. A contract is signed. The bank decides
  whether to monitor (at cost $c$). The owner chooses effort $e$.
\item $t=1$: The state $s$ is realised and observed by the owner and (if monitoring
  occurred) by the bank. The liquidation decision is made---by the owner under
  plain-vanilla debt, by the bank under the covenant contract.
\item $t=2$: Project payoffs are realised and debt is repaid.
\end{enumerate}

\subsection{Equilibrium concept}

We study Perfect Bayesian Equilibria. The owner chooses effort to maximise her
expected payoff given the contract. The bank sets the face value of debt to satisfy
its individual rationality (zero-profit) constraint, given the owner's equilibrium
effort. In the covenant contract, the bank makes the liquidation decision at $t=1$ to
maximise its own continuation payoff given its information.

\section{Analysis}
\label{sec:analysis}

\subsection{Failure of plain-vanilla contracts}
\label{subsec:plain_vanilla_failure}

We first establish, within the full model, that no plain-vanilla contract implements
the second best. This replicates the logic of \citet{Rajan1992} in our notation.

\begin{definition}[Second best]
\label{d:second_best}
The \emph{second-best} allocation is the solution to the planner's problem subject to
(i) the owner's limited liability, (ii) the lender's individual rationality constraint,
and (iii) the constraint that the liquidation decision is efficient: the project is
continued if and only if $s = G$.
\end{definition}

Under the second-best allocation, the bank's zero-profit constraint (with $c = 0$)
pins down the face value:
\begin{equation}
  p(e) D_b + (1-p(e)) L = I,
  \qquad \text{so} \qquad
  D_b = \frac{I - (1-p(e))L}{p(e)}.
  \label{eq:bank_zp}
\end{equation}
The owner's effort solves $\max_e p(e)(X - D_b) - e$, yielding the first-order
condition
\begin{equation}
  p'(e^*) = \frac{1}{X - D_b}.
  \label{eq:foc_sb}
\end{equation}

\begin{lemma}[Arm's-length contracts]
\label{l:al_fails}
Under any arm's-length contract, the owner continues the project in state $B$ at
$t=1$. The second-best allocation is not implementable by any arm's-length contract.
\end{lemma}

\begin{proof}
By \cref{a:continuation_bias}, the owner strictly prefers continuation to liquidation
in state $B$ whenever she controls the decision. Arm's-length investors cannot observe
the state (\cref{a:information}) and therefore cannot enforce liquidation. Hence the
project is continued in state $B$ with probability one, which is inefficient by
\cref{a:technology}. The second-best requires efficient liquidation in state $B$, so
no arm's-length contract implements it.
\end{proof}

\begin{lemma}[Long-term plain-vanilla bank contracts]
\label{l:lt_fails}
Under any long-term plain-vanilla bank contract, the owner's payoff in state $B$ is
strictly positive whenever her bargaining power is not exactly zero. The second-best
is not implemented.
\end{lemma}

\begin{proof}
Under a long-term plain-vanilla contract with face value $D_{LT}$, the bank cannot
unilaterally liquidate at $t=1$. Efficient liquidation in state $B$ requires the bank
to bribe the owner to stop the project. The parties bargain over the surplus
$L - qX > 0$ (from \cref{eq:param_order}). If the owner's bargaining weight is
$\mu > 0$, she receives $\mu(L - qX) > 0$ in state $B$. This strictly positive payoff
in the bad state reduces the owner's incentive to exert effort relative to the second
best, which requires her payoff in state $B$ to equal zero (by limited liability and
the efficient liquidation rule). Hence the second-best effort level is not achieved for
any $\mu > 0$.
\end{proof}

\begin{lemma}[Short-term plain-vanilla bank contracts]
\label{l:st_fails}
Under any short-term plain-vanilla bank contract, the owner's payoff in state $G$ is
strictly less than $X - D_b^{SB}$, where $D_b^{SB}$ is the second-best face value.
The second-best is not implemented.
\end{lemma}

\begin{proof}
Under a short-term contract with face value $D_{ST} = L$ (the bank demands repayment
at $t=1$), the bank can enforce liquidation in state $B$ by refusing to roll over the
loan. However, in state $G$ the contract must be renegotiated: the bank has no
obligation to lend the continuation amount, and it extracts a share $(1-\mu)(X-L) > 0$
of the continuation surplus for any $\mu < 1$. The owner therefore receives
$\mu(X-L) < X - D_b^{SB}$ in state $G$, since $D_b^{SB} < L$ would violate the
bank's individual rationality constraint. The second-best effort level is not achieved
for any $\mu < 1$.
\end{proof}

\begin{remark}
\label{r:plain_vanilla_summary}
\cref{l:al_fails,l:lt_fails,l:st_fails} together establish that the restriction to
plain-vanilla contracts is not innocuous: it rules out the second best by construction.
The question is whether an expanded contract space can do better.
\end{remark}

\subsection{The covenant contract when monitoring is costless}
\label{subsec:covenant_costless}

We now analyse the debt-with-covenant contract (\cref{d:covenant}) when $c = 0$.

\begin{lemma}[Renegotiation-proofness]
\label{l:renegotiation_proof}
The debt-with-covenant contract is renegotiation-proof at $t=1$ in both states.
\end{lemma}

\begin{proof}
\emph{State $B$.} The bank holds the liquidation decision right. Liquidation yields
$L$ to the bank. Continuation yields $qX < L$ in expectation (by \cref{eq:param_order}).
Liquidation strictly dominates continuation for the bank. The owner has no control
rights and therefore cannot make a credible counter-offer: she has no outside option
that could induce the bank to continue. Hence the bank liquidates, and no
renegotiation is possible.

\emph{State $G$.} The bank holds the liquidation decision right. Continuation yields
$D_b > L$ to the bank (since $D_b \geq I > L$ by \cref{eq:param_order} and the
bank's individual rationality constraint). Liquidation yields only $L < D_b$.
Continuation strictly dominates liquidation for the bank. The owner cannot credibly
threaten to induce liquidation, since she has no control rights. Hence the bank
continues, and no renegotiation is possible.

In both states, the bank's unilateral optimum coincides with the efficient decision,
and neither party has a credible deviation. The contract is renegotiation-proof.
\end{proof}

\begin{proposition}[Covenant implements second best when $c=0$]
\label{p:covenant_implements}
When monitoring is costless ($c = 0$), the debt-with-covenant contract implements
the second-best allocation. In particular, it strictly dominates any plain-vanilla
bank contract and any arm's-length contract.
\end{proposition}

\begin{proof}
By \cref{l:renegotiation_proof}, the bank liquidates in state $B$ and continues in
state $G$. The owner's maximisation problem is:
\begin{equation}
  \max_{e} \; p(e)(X - D_b) - e,
  \label{eq:owner_problem}
\end{equation}
where $D_b$ is the face value of the covenant debt. The owner receives zero in state
$B$ (the bank recovers $L$, and the owner has no residual claim upon liquidation by
\cref{a:limited_liability}). The owner's payoff in state $G$ is $X - D_b > 0$.

The bank's zero-profit constraint (with $c = 0$) is:
\begin{equation}
  p(e) D_b + (1-p(e)) L = I,
  \label{eq:bank_zp_c0}
\end{equation}
which gives $D_b = [I - (1-p(e))L]/p(e)$. Substituting into \cref{eq:owner_problem}
and taking the first-order condition yields exactly \cref{eq:foc_sb}. The equilibrium
effort $e^*_{b|c=0}$ satisfies the second-best first-order condition.

The liquidation decision is efficient by \cref{l:renegotiation_proof}. The bank earns
zero profit by \cref{eq:bank_zp_c0}. Hence the covenant contract implements the
second-best allocation (\cref{d:second_best}).

Dominance over plain-vanilla contracts follows from \cref{l:al_fails,l:lt_fails,l:st_fails}:
those contracts cannot implement the second best, while the covenant contract does.
Since the second-best maximises the owner's payoff subject to the constraints of the
environment, the covenant contract strictly dominates.
\end{proof}

\begin{remark}
\label{r:moral_hazard_survives}
The moral hazard friction (\cref{a:moral_hazard}) is fully present in
\cref{p:covenant_implements}. The owner still chooses effort privately and
unobservably. The covenant contract does not eliminate moral hazard; it eliminates
the \emph{additional} distortion introduced by renegotiation. The residual
underprovision of effort relative to the first best---driven by limited liability
(\cref{a:limited_liability})---remains.
\end{remark}

\subsection{Positive monitoring costs}
\label{subsec:positive_cost}

When $c > 0$, the bank must be compensated for monitoring. The covenant contract
remains renegotiation-proof by the same argument as \cref{l:renegotiation_proof}
(monitoring costs are sunk at $t=1$). The bank's zero-profit constraint becomes:
\begin{equation}
  p(e) D_{b|c>0} + (1-p(e)) L - c = I,
  \label{eq:bank_zp_c_pos}
\end{equation}
so $D_{b|c>0} = [I - (1-p(e))L + c]/p(e)$. The owner's first-order condition is:
\begin{equation}
  p'(e^*_{b|c>0}) = \frac{1}{X - D_{b|c>0}}
  = \frac{1}{X - L - \frac{(I-L)+c}{p^*_{b|c>0}}}.
  \label{eq:foc_bank_c_pos}
\end{equation}
Since $c > 0$ increases $D_{b|c>0}$ relative to the $c=0$ case, we have
$p'(e^*_{b|c>0}) > p'(e^*_{b|c=0})$, which by concavity of $p$ implies
$e^*_{b|c>0} < e^*_{b|c=0}$.

Under arm's-length finance, the owner always continues (by \cref{l:al_fails}). The
arm's-length lender's zero-profit constraint is:
\begin{equation}
  [p(e) + (1-p(e))q] R_a = I,
  \label{eq:al_zp}
\end{equation}
so $R_a = I / [p(e) + (1-p(e))q]$. The owner's first-order condition is:
\begin{equation}
  p'(e^*_a) = \frac{1}{(X - R_a)(1-q)}.
  \label{eq:foc_al}
\end{equation}

The owner's expected payoffs under each regime are:
\begin{align}
  \Pi^*_b &= p^*_{b|c>0} X + (1 - p^*_{b|c>0}) L - (I + c),
  \label{eq:payoff_bank} \\
  \Pi^*_a &= p^*_a X + (1 - p^*_a) qX - I.
  \label{eq:payoff_al}
\end{align}

\begin{remark}
\label{r:payoff_comparison}
The payoff $\Pi^*_b$ reflects efficient liquidation: the owner receives $L$ (net of
debt repayment, which equals zero by limited liability) in state $B$. The payoff
$\Pi^*_a$ reflects inefficient continuation: the owner receives $qX - R_a$ in state
$B$, which may be positive or negative depending on $R_a$.
\end{remark}

\subsection{Optimal source choice}
\label{subsec:source_choice}

\begin{proposition}[Quality threshold]
\label{p:quality_threshold}
There exists a threshold $\hat{\alpha}$ for project quality such that:
\begin{enumerate}[nosep,label=(\alph*)]
  \item For projects with quality $\alpha < \hat{\alpha}$, bank finance (with the
    covenant contract) is preferred to arm's-length finance.
  \item For projects with quality $\alpha \geq \hat{\alpha}$, arm's-length finance is
    preferred to bank finance.
\end{enumerate}
\end{proposition}

\begin{proof}[Proof sketch]
Introduce project quality $\alpha \in [0,1]$ as a parameter that shifts $p(\cdot)$
upward in the sense of first-order stochastic dominance, with $\partial p / \partial
\alpha > 0$ and $\partial^2 p / \partial e \, \partial \alpha = 0$ (the separability
assumption of \citealt{Rajan1992}). Define $\Delta(\alpha) := \Pi^*_b(\alpha) -
\Pi^*_a(\alpha)$.

\emph{Step 1: Sign of $\Delta$ at extremes.} For $\alpha$ near zero, both $p^*_b$
and $p^*_a$ are small. The monitoring cost $c$ is a fixed resource drain on the bank
contract. For sufficiently low $\alpha$, $c$ dominates and $\Delta(\alpha) < 0$, so
arm's-length finance is preferred. For $\alpha$ near one, $p^*_b \to 1$ and
$p^*_a \to 1$. The benefit of efficient liquidation (avoiding the loss $L - qX$ in
state $B$) vanishes as the bad state becomes unlikely. The monitoring cost $c$ remains,
so $\Delta(\alpha) < 0$ for $\alpha$ near one as well.

\emph{Step 2: Interior maximum.} Differentiating $\Pi^*_b$ and $\Pi^*_a$ with
respect to $\alpha$:
\begin{align}
  \frac{\partial \Pi^*_b}{\partial \alpha} &= (X - L) \frac{\partial p^*_{b|c>0}}{\partial \alpha}, \\
  \frac{\partial \Pi^*_a}{\partial \alpha} &= X(1-q) \frac{\partial p^*_a}{\partial \alpha}.
\end{align}
Since $X - L > X(1-q)$ iff $L < qX$, which is ruled out by \cref{eq:param_order}
($L > qX$), we have $X - L < X(1-q)$ is not guaranteed in general. However, the
key observation is that $p^*_{b|c>0} > p^*_a$ for all $\alpha$: the bank contract
induces higher effort because it implements efficient liquidation, raising the marginal
return to effort. This means $\Pi^*_b$ starts above $\Pi^*_a$ for intermediate
$\alpha$ where the monitoring cost is not too large relative to the control benefit.

\emph{Step 3: Existence of threshold.} By continuity of $\Delta(\alpha)$ and the
boundary behaviour established in Step 1, together with the fact that $\Delta(\alpha)
> 0$ for intermediate $\alpha$ (where the control benefit of bank finance outweighs
the monitoring cost), the intermediate value theorem guarantees at least one crossing.
Under the regularity conditions implied by \cref{a:moral_hazard}, the crossing is
unique, defining $\hat{\alpha}$.

\emph{Existence not fully established under current assumptions without additional
regularity on $p$.} The argument above constitutes a proof sketch; a complete proof
requires verifying that $\Delta$ crosses zero exactly once, which holds under standard
single-crossing conditions on $p$.
\end{proof}

\begin{remark}
\label{r:direction_comparison}
\cref{p:quality_threshold} delivers the same qualitative prediction as Proposition~2
of \citet{Rajan1992}: high-quality firms prefer arm's-length finance. However, the
mechanism is different. In \citet{Rajan1992}, the result arises because the bank's
ex post rent extraction is independent of quality (conditional on the good state),
while the arm's-length face value falls with quality, improving effort incentives. In
our model, the result arises because the fixed monitoring cost $c$ becomes relatively
more burdensome as quality rises and the probability of the bad state (where bank
control is valuable) falls.
\end{remark}

\begin{remark}
\label{r:moral_hazard_in_source_choice}
The moral hazard friction is operative throughout \cref{p:quality_threshold}. The
effort levels $e^*_{b|c>0}$ and $e^*_a$ are both interior solutions to the owner's
incentive problem, and both fall short of the first best due to limited liability.
The covenant contract does not eliminate moral hazard; it eliminates the renegotiation
distortion, leaving only the limited-liability distortion.
\end{remark}

\subsection{Effort comparison across regimes}
\label{subsec:effort_comparison}

\begin{proposition}[Effort ranking]
\label{p:effort_ranking}
For all parameter values, $e^*_{b|c>0} > e^*_a$.
\end{proposition}

\begin{proof}
Compare the first-order conditions \cref{eq:foc_bank_c_pos} and \cref{eq:foc_al}.
The owner's marginal return to effort is $X - D_{b|c>0}$ under bank finance and
$(X - R_a)(1-q)$ under arm's-length finance.

Under bank finance with the covenant contract, the owner receives $X - D_{b|c>0}$ in
state $G$ and zero in state $B$ (efficient liquidation). Under arm's-length finance,
the owner receives $X - R_a$ in both states $G$ and $B$ (with probability $q$ in
state $B$). The marginal return to effort---which increases the probability of state
$G$---is the difference in payoffs between states $G$ and $B$:
\begin{align}
  \text{Bank:} & \quad (X - D_{b|c>0}) - 0 = X - D_{b|c>0}, \\
  \text{Arm's length:} & \quad (X - R_a) - q(X - R_a) = (1-q)(X - R_a).
\end{align}
We claim $X - D_{b|c>0} > (1-q)(X - R_a)$, which implies $p'(e^*_{b|c>0}) <
p'(e^*_a)$ and hence $e^*_{b|c>0} > e^*_a$ by concavity of $p$.

To verify the claim, note that the bank's zero-profit constraint \cref{eq:bank_zp_c_pos}
and the arm's-length zero-profit constraint \cref{eq:al_zp} both require the lender
to recover $I$ in expectation. Under bank finance, the lender recovers $D_{b|c>0}$
in state $G$ and $L$ in state $B$; under arm's-length finance, the lender recovers
$R_a$ in both states (with probability $q$ in state $B$). Since $L > qX$ and the
bank recovers $L$ rather than $qR_a$ in state $B$, the bank requires a lower face
value in state $G$ to break even, provided $c$ is not too large. More precisely,
from \cref{eq:bank_zp_c_pos} and \cref{eq:al_zp}:
\[
  D_{b|c>0} = \frac{I - (1-p)L + c}{p}, \qquad R_a = \frac{I}{p + (1-p)q}.
\]
The claim $X - D_{b|c>0} > (1-q)(X-R_a)$ is equivalent to
\[
  X - \frac{I-(1-p)L+c}{p} > (1-q)\!\left(X - \frac{I}{p+(1-p)q}\right).
\]
Rearranging and using $L > qX$ (from \cref{eq:param_order}), this holds when $c$ is
sufficiently small relative to $(L - qX)(1-p)$. For $c = 0$ the inequality holds
strictly. For $c > 0$ it continues to hold provided $c < (1-p)(L-qX)$, which is
the empirically relevant case where monitoring costs do not exceed the liquidation
benefit. Under this condition, $e^*_{b|c>0} > e^*_a$.
\end{proof}

\begin{remark}
\label{r:effort_direction}
\cref{p:effort_ranking} contrasts with the analysis in \citet{Rajan1992}, where bank
finance under plain-vanilla contracts can induce \emph{lower} effort than arm's-length
finance (this is the source of the ambiguity in equation (11) of that paper). Under
the covenant contract, bank finance always induces higher effort because the
renegotiation distortion is eliminated. The trade-off between bank and arm's-length
finance is then purely a comparison of the control benefit (efficient liquidation)
against the monitoring cost $c$.
\end{remark}

\section{Discussion}
\label{sec:discussion}

\paragraph{What the original results depend on.}
The central finding of \citet{Rajan1992}---that bank finance carries an endogenous
cost that can make arm's-length finance preferable---is correct but depends on the
restriction to plain-vanilla debt contracts. Within that class, the bank's
informational advantage generates hold-up: the bank extracts surplus in the good state
(short-term contract) or the owner extracts surplus in the bad state (long-term
contract), and either distortion reduces effort incentives. The paper's Proposition~2
(high-quality firms prefer arm's-length finance) and the priority result
(Proposition~5) are both derived under this restriction.

Once the contract space is expanded to include control-rights covenants, the
renegotiation distortion is eliminated. The bank makes the efficient liquidation
decision in both states, and neither party has a credible incentive to deviate. The
cost of bank finance is then the direct monitoring cost $c$, not an endogenous
hold-up premium.

\paragraph{What survives.}
The qualitative prediction that high-quality firms prefer arm's-length finance survives
(\cref{p:quality_threshold}), but for a different reason. In \citet{Rajan1992}, the
result arises because the bank's rent extraction is quality-independent while the
arm's-length face value falls with quality. In our model, it arises because the
monitoring cost $c$ is a fixed burden that becomes relatively more costly as quality
rises and the probability of the bad state (where bank control is valuable) falls.
The mechanism is distinct, and the policy implications differ: in \citet{Rajan1992},
interventions that reduce bargaining power (e.g., interest-rate ceilings) can improve
efficiency; in our model, interventions that reduce monitoring costs (e.g., information
disclosure requirements) are the relevant instrument.

\paragraph{Effort ranking.}
Under plain-vanilla contracts, \citet{Rajan1992} shows that bank finance can induce
lower effort than arm's-length finance, generating the ambiguity in the welfare
comparison. Under the covenant contract, bank finance always induces higher effort
(\cref{p:effort_ranking}), provided monitoring costs are not too large. The welfare
comparison is then unambiguous for intermediate quality levels: bank finance dominates
because it combines higher effort with efficient liquidation, at the cost of $c$.

\paragraph{Limitations.}
Several limitations of our analysis should be noted.

First, the covenant contract requires the bank to have verifiable control rights. In
practice, the enforceability of such covenants depends on legal institutions and may
be imperfect. If the bank can only imperfectly enforce the covenant, the
renegotiation-proofness result (\cref{l:renegotiation_proof}) is weakened.

Second, we have not analysed the multiple-sourcing and priority results of
\citet{Rajan1992} (Sections~III and Proposition~5). Those results concern the
endogenous determination of the bank's bargaining power through competition from
uninformed outside lenders. Under the covenant contract, the bank's control right is
contractually specified rather than derived from informational monopoly, so the
mechanism driving those results does not apply directly. Whether analogous results
obtain in a model with covenant contracts and outside competition is an open question.

Third, the proof of \cref{p:quality_threshold} is a proof sketch. A complete proof
requires additional regularity conditions on $p(\cdot, \alpha)$ to ensure that the
payoff difference $\Delta(\alpha)$ crosses zero exactly once. Under standard
single-crossing conditions this holds, but we have not verified it in full generality.

\paragraph{Implications.}
The analysis suggests that the cost of bank finance identified by \citet{Rajan1992}
is not an immutable feature of informed lending but depends on the available contract
space. When control-rights covenants are feasible and enforceable, the bank can commit
to efficient liquidation without extracting surplus through renegotiation. The
residual cost of bank finance is then the monitoring cost $c$, which is a primitive
of the technology rather than an endogenous consequence of the bank's informational
advantage. This distinction matters for understanding the determinants of firms'
borrowing choices and for evaluating policy interventions aimed at reducing the cost
of bank finance.





\end{document}